\documentclass[aspectratio=169]{beamer}

\usetheme{Madrid}
\usecolortheme{default}

\usepackage[utf8]{inputenc}
\usepackage{listings}
\usepackage{xcolor}
\usepackage{amsmath}

\definecolor{codegray}{RGB}{245,245,245}
\definecolor{codeblue}{RGB}{17,24,39}
\definecolor{codegreen}{RGB}{30,120,60}

\lstset{
    backgroundcolor=\color{codegray},
    basicstyle=\ttfamily\scriptsize,
    keywordstyle=\color{codeblue},
    commentstyle=\color{codegreen},
    stringstyle=\color{red!70!black},
    breaklines=true,
    frame=single,
    showstringspaces=false
}

\title{Retrieval-Augmented Generation (RAG)}
%\subtitle{A Beginner Tutorial with Python}
\author{Dr. Sisay Chala}
%\date{\today}

\begin{document}

%------------------------------------------------
\begin{frame}
    \titlepage
\end{frame}

%------------------------------------------------
\begin{frame}{Learning Objectives}

By the end of this presentation, you will understand:

\begin{itemize}[<+->]
    \item What Retrieval-Augmented Generation (RAG) is
    \item Why RAG is useful
    \item What document chunking means
    \item What embeddings are
    \item What FAISS does
    \item How retrieval works
    \item How retrieved information is given to an LLM
    \item How to build a simple RAG pipeline in Python
\end{itemize}

\end{frame}

%------------------------------------------------
\begin{frame}{What is RAG?}

\textbf{RAG = Retrieval-Augmented Generation}

Instead of asking an LLM to answer only from its training knowledge:

\begin{center}
\Large
Question $\rightarrow$ LLM $\rightarrow$ Answer
\end{center}

\vspace{0.5cm}

RAG first retrieves relevant information from external documents:

\begin{center}
\Large
Question
$\rightarrow$
Retrieve relevant information
$\rightarrow$
LLM
$\rightarrow$
Answer
\end{center}

\end{frame}

%------------------------------------------------
\begin{frame}{Why Use RAG?}

Imagine we have a company handbook:

\begin{block}{Company Policy}
Employees receive 30 days of paid vacation per year.
\end{block}

A user asks:

\begin{quote}
How many vacation days do employees receive?
\end{quote}

A general-purpose LLM may not know this private company policy.

\vspace{0.3cm}

RAG retrieves the relevant passage and provides it to the LLM.

\end{frame}

%------------------------------------------------
\begin{frame}[shrink]{Basic RAG Architecture}

\begin{center}

\begin{tabular}{c}
Documents \\
$\downarrow$ \\
Chunking \\
$\downarrow$ \\
Embeddings \\
$\downarrow$ \\
Vector Database / FAISS \\
$\downarrow$ \\
User Question \\
$\downarrow$ \\
Similarity Search \\
$\downarrow$ \\
Relevant Chunks \\
$\downarrow$ \\
LLM \\
$\downarrow$ \\
Answer
\end{tabular}

\end{center}

\end{frame}

%------------------------------------------------
\begin{frame}[fragile]{Step 1: Load Documents}

Create a simple document:

\begin{lstlisting}[language={},caption={documents.txt}]
RAG stands for Retrieval-Augmented Generation.

RAG systems retrieve relevant information from a
knowledge base and provide that information to a
language model as context.

The main stages of a RAG system are document loading,
chunking, embedding, retrieval, and generation.
\end{lstlisting}

Load it with Python:

\begin{lstlisting}[language=Python]
with open("documents.txt", "r",
          encoding="utf-8") as f:
    text = f.read()
\end{lstlisting}

\end{frame}

%------------------------------------------------
\begin{frame}[fragile]{Step 2: Chunking}

Large documents are divided into smaller pieces called
\textbf{chunks}.

\begin{lstlisting}[language=Python]
def chunk_text(text, chunk_size=300):
    words = text.split()

    chunks = []

    for i in range(0, len(words), chunk_size):
        chunk = " ".join(
            words[i:i + chunk_size]
        )
        chunks.append(chunk)

    return chunks
\end{lstlisting}

\end{frame}

%------------------------------------------------
\begin{frame}[fragile]{What is chunk\_size?}

In our simple implementation:

\begin{lstlisting}[language=Python]
def chunk_text(text, chunk_size=300):
\end{lstlisting}

\texttt{chunk\_size=300} means:

\begin{center}
\Large
\textbf{up to 300 words per chunk}
\end{center}

For example, with a 750-word document:

\begin{itemize}
    \item Chunk 1: words 1--300
    \item Chunk 2: words 301--600
    \item Chunk 3: words 601--750
\end{itemize}

\end{frame}

%------------------------------------------------
\begin{frame}{Why Does Chunk Size Matter?}

\textbf{Chunks that are too small}

\begin{itemize}
    \item May lose important context
    \item Retrieved information may be incomplete
\end{itemize}

\vspace{0.3cm}

\textbf{Chunks that are too large}

\begin{itemize}
    \item May contain irrelevant information
    \item Retrieval becomes less precise
    \item More context is sent to the LLM
\end{itemize}

\vspace{0.3cm}

In real systems, chunks often use:

\begin{itemize}
    \item Paragraphs
    \item Sections
    \item Sentences
    \item Token-based limits
    \item Overlapping chunks
\end{itemize}

\end{frame}

%------------------------------------------------
\begin{frame}{Step 3: Embeddings}

An \textbf{embedding} converts text into a numerical vector.

For example:

\begin{center}
\texttt{"RAG retrieves documents"}
\end{center}

becomes something like:

\[
[0.12,\;-0.44,\;0.81,\;0.03,\;\ldots]
\]

The vector represents semantic information about the text.

\vspace{0.3cm}

Similar pieces of text tend to have similar vectors.

\end{frame}

%------------------------------------------------
\begin{frame}[fragile, shrink]{Creating Embeddings with Python}

Install Sentence Transformers:

\begin{lstlisting}[language=bash]
pip install sentence-transformers
\end{lstlisting}

Load an embedding model:

\begin{lstlisting}[language=Python]
from sentence_transformers import SentenceTransformer

model = SentenceTransformer(
    "all-MiniLM-L6-v2"
)
\end{lstlisting}

Create embeddings:

\begin{lstlisting}[language=Python]
embeddings = model.encode(chunks)

print(embeddings.shape)
\end{lstlisting}

For example:
\[
(2,384)
\]
means 2 chunks with 384-dimensional vectors.

\end{frame}

%------------------------------------------------
\begin{frame}{Step 4: Indexing using FAISS}

\textbf{FAISS = Facebook AI Similarity Search}

FAISS is a library developed by Meta for
\textbf{efficient similarity search over vectors}.

In RAG, FAISS can store document embeddings and
quickly find vectors similar to a user's question.

\vspace{0.4cm}

Think of FAISS as:

\begin{center}
\Large
\textbf{a search engine for embeddings}
\end{center}

\end{frame}

%------------------------------------------------
\begin{frame}{FAISS in the RAG Pipeline}

\begin{center}

\begin{tabular}{c}
Document chunks \\
$\downarrow$ \\
Embedding model \\
$\downarrow$ \\
Vectors \\
$\downarrow$ \\
\textbf{FAISS} \\
$\downarrow$ \\
Similar vectors / relevant chunks
\end{tabular}

\end{center}

FAISS does \textbf{not} generate answers.

It performs similarity search.

\end{frame}

%------------------------------------------------
\begin{frame}[fragile]{Creating a FAISS Index}

Install FAISS:

\begin{lstlisting}[language=bash]
pip install faiss-cpu
\end{lstlisting}

Create an index:

\begin{lstlisting}[language=Python]
import faiss

dimension = embeddings.shape[1]

index = faiss.IndexFlatL2(dimension)
\end{lstlisting}

Add document embeddings:

\begin{lstlisting}[language=Python]
index.add(embeddings)
\end{lstlisting}

\end{frame}

%------------------------------------------------
\begin{frame}[fragile]{What Does IndexFlatL2 Mean?}

Our example uses:

\begin{lstlisting}[language=Python]
faiss.IndexFlatL2(dimension)
\end{lstlisting}

\texttt{L2} refers to Euclidean distance.

Conceptually:

\[
d(x,y) =
\sqrt{\sum_i (x_i-y_i)^2}
\]

Smaller distance means the vectors are closer.

\vspace{0.3cm}

FAISS provides many other index types optimized for
different datasets and performance requirements.

\end{frame}

%------------------------------------------------
\begin{frame}[fragile]{Step 5: Retrieval}

Suppose the user asks:

\begin{quote}
What are the main stages of RAG?
\end{quote}

First, embed the question:

\begin{lstlisting}[language=Python]
question_embedding = model.encode(
    [question]
)
\end{lstlisting}

Then search FAISS:

\begin{lstlisting}[language=Python]
distances, indices = index.search(
    question_embedding,
    k=2
)
\end{lstlisting}

\texttt{k=2} means:

\begin{center}
\textbf{retrieve the 2 most similar chunks}
\end{center}

\end{frame}

%------------------------------------------------
\begin{frame}[fragile]{Retrieving the Actual Text}

FAISS returns the indices of the matching vectors.

We use those indices to retrieve the original chunks:

\begin{lstlisting}[language=Python]
retrieved_chunks = [
    chunks[i]
    for i in indices[0]
]
\end{lstlisting}

Now we have the information that is
most relevant to the question.

\end{frame}

%------------------------------------------------
\begin{frame}[fragile]{Step 6: Give Context to the LLM}

Combine the retrieved chunks:

\begin{lstlisting}[language=Python]
context = "\n\n".join(
    retrieved_chunks
)
\end{lstlisting}

Build a prompt:

\begin{lstlisting}[language=Python]
prompt = f"""
Answer the question using only the context below.

Context:
{context}

Question:
{question}

If the answer cannot be found in the context, say:
"I don't know based on the provided documents."
"""
\end{lstlisting}

\end{frame}

%------------------------------------------------
\begin{frame}[fragile]{Step 7: Generate the Answer}

Use an LLM:

\begin{lstlisting}[language=Python]
from openai import OpenAI

client = OpenAI()

response = client.responses.create(
    model="gpt-5.6",
    input=prompt
)

print(response.output_text)
\end{lstlisting}

% The LLM receives:

% \[
% \text{Question}
% +
% \text{Retrieved Context}
% \rightarrow
% \text{Answer}
% \]

\end{frame}

%------------------------------------------------
\begin{frame}{Complete RAG Pipeline}

\begin{center}
\Large

\textbf{Documents}

$\downarrow$

\textbf{Chunking}

$\downarrow$

\textbf{Embedding Model}

$\downarrow$

\textbf{FAISS}

$\downarrow$

\textbf{Question}

$\downarrow$

\textbf{Question Embedding}

$\downarrow$

\textbf{Similarity Search}

$\downarrow$

\textbf{Relevant Chunks}

$\downarrow$

\textbf{LLM}

$\downarrow$

\textbf{Answer}

\end{center}

\end{frame}

%------------------------------------------------
\begin{frame}{The Three Main Components}

A simple RAG system combines three important components.

\begin{block}{1. Embedding Model}
Converts text into vectors.
\end{block}

\begin{block}{2. Vector Search}
Finds document vectors similar to the question.
FAISS can perform this task.
\end{block}

\begin{block}{3. LLM}
Uses the retrieved context to generate a natural-language answer.
\end{block}

\end{frame}

%------------------------------------------------
\begin{frame}{Traditional LLM vs RAG}

\begin{columns}

\column{0.48\textwidth}

\textbf{Traditional LLM}

\begin{center}
Question

$\downarrow$

LLM

$\downarrow$

Answer
\end{center}

The model relies primarily on
its existing knowledge.

\pause 

\column{0.48\textwidth}

\textbf{RAG}

\begin{center}
Question

$\downarrow$

Retrieve documents

$\downarrow$

LLM

$\downarrow$

Answer
\end{center}

The model receives relevant
external context.

\end{columns}

\end{frame}

%------------------------------------------------
\begin{frame}{Where Can RAG Be Used?}

RAG is useful when an LLM needs access to
external or private information.

\begin{itemize}
    \item Company knowledge bases
    \item Internal documentation
    \item Customer-support systems
    \item Research assistants
    \item Product manuals
    \item Technical documentation
    \item Question answering over PDFs
    \item Internal wikis
\end{itemize}

\end{frame}

% %------------------------------------------------
% \begin{frame}[fragile]{Simple RAG: Complete Example}

% \begin{lstlisting}[language=Python]
% import faiss
% from sentence_transformers import SentenceTransformer
% from openai import OpenAI

% # Load document
% with open("documents.txt", "r",
%           encoding="utf-8") as f:
%     text = f.read()

% # Chunk document
% def chunk_text(text, chunk_size=100):
%     words = text.split()
%     return [
%         " ".join(words[i:i + chunk_size])
%         for i in range(0, len(words), chunk_size)
%     ]

% chunks = chunk_text(text)

% # Create embeddings
% model = SentenceTransformer(
%     "all-MiniLM-L6-v2"
% )

% embeddings = model.encode(chunks)

% # Create FAISS index
% dimension = embeddings.shape[1]
% index = faiss.IndexFlatL2(dimension)
% index.add(embeddings)
% \end{lstlisting}

% \end{frame}

% %------------------------------------------------
% \begin{frame}[fragile]{Complete Example: Retrieval}

% \begin{lstlisting}[language=Python]
% question = "What are the main stages of RAG?"

% # Embed question
% question_embedding = model.encode(
%     [question]
% )

% # Search
% distances, indices = index.search(
%     question_embedding,
%     k=2
% )

% # Retrieve text
% retrieved_chunks = [
%     chunks[i]
%     for i in indices[0]
% ]

% context = "\n\n".join(
%     retrieved_chunks
% )
% \end{lstlisting}

% \end{frame}

% %------------------------------------------------
% \begin{frame}[fragile]{Complete Example: Generation}

% \begin{lstlisting}[language=Python]
% prompt = f"""
% Answer the question using only
% the context below.

% Context:
% {context}

% Question:
% {question}

% If the answer cannot be found in
% the context, say:
% "I don't know based on the
% provided documents."
% """

% client = OpenAI()

% response = client.responses.create(
%     model="gpt-5.6",
%     input=prompt
% )

% print(response.output_text)
% \end{lstlisting}

% \end{frame}

%------------------------------------------------
\begin{frame}{Limitations of our RAG Example}

%Our implementation is intentionally simple.

A production RAG system needs to consider:

\begin{itemize}
    \item Better document parsing
    \item Better chunking strategies
    \item \textbf{Chunk overlap}
    \item \textbf{Metadata}
    \item Better vector indexes
    \item Hybrid search
    \item Reranking
    \item Retrieval evaluation
    \item Answer evaluation
    \item Source citations
\end{itemize}

\end{frame}

%------------------------------------------------
\begin{frame}{Chunk Overlap}

Instead of completely independent chunks:

\begin{center}
Chunk 1: 1--300

Chunk 2: 301--600
\end{center}

we can use overlap:

\begin{center}
Chunk 1: 1--300

Chunk 2: 250--550

Chunk 3: 500--800
\end{center}

Overlap helps preserve information that occurs
near chunk boundaries.

\end{frame}

% %------------------------------------------------
% \begin{frame}{Advanced Retrieval}

% A production system may use:

% \begin{itemize}
%     \item Semantic vector search
%     \item Keyword search
%     \item Metadata filtering
%     \item Hybrid search
%     \item Reranking
% \end{itemize}

% A common pipeline is:

% \begin{center}
% Question
% $\rightarrow$
% Retrieve 20 candidates
% $\rightarrow$
% Reranker
% $\rightarrow$
% Top 5 chunks
% $\rightarrow$
% LLM
% \end{center}

% \end{frame}

%------------------------------------------------
\begin{frame}[fragile]{Metadata}

Instead of storing only the text:

\begin{lstlisting}[language=Python]
chunk
\end{lstlisting}

we can store additional information:

\begin{lstlisting}[language=Python]
{
    "text": "...",
    "document": "handbook.pdf",
    "page": 42,
    "section": "Vacation Policy"
}
\end{lstlisting}

This makes it possible to provide source information
with the generated answer.

\end{frame}

\begin{fragile}{Example Implementation}

You can find the example implementation on GitHub: 

\url{https://github.com/AI-Cloud-AWS/rag-example}

In order to execute this example, you need:
    \begin{itemize}
        \item 
        \item OpenAI key 
    \end{itemize}
% %------------------------------------------------
% \begin{frame}{Key Concepts to Learn Next}

% A useful learning sequence is:

% \begin{enumerate}
%     \item Embeddings
%     \item Vector similarity
%     \item Chunking strategies
%     \item Vector databases
%     \item Metadata filtering
%     \item Hybrid search
%     \item Reranking
%     \item Context windows
%     \item RAG evaluation
%     \item Advanced RAG
% \end{enumerate}

% \end{frame}

%------------------------------------------------
\begin{frame}{Recap}

\begin{center}

\Huge
RAG is not a special type of LLM.

\vspace{0.5cm}

\Large
It is a pipeline combining

\vspace{0.3cm}

\textbf{Information Retrieval}

\vspace{0.2cm}

+

\vspace{0.2cm}

\textbf{Language Generation}

\end{center}

\end{frame}

% %------------------------------------------------
% \begin{frame}{Summary}

% \begin{itemize}
%     \item Documents are divided into chunks.
%     \item Chunks are converted into embeddings.
%     \item FAISS stores and searches the embeddings.
%     \item A user's question is also converted into an embedding.
%     \item Similar document chunks are retrieved.
%     \item Retrieved chunks are added to the LLM prompt.
%     \item The LLM generates an answer using the retrieved context.
% \end{itemize}

% \vspace{0.5cm}

% \begin{center}
% \Large
% \textbf{Question $\rightarrow$ Retrieve $\rightarrow$ Generate}
% \end{center}

% \end{frame}

%------------------------------------------------
\begin{frame}
    \centering
    \Huge
    Questions?
\end{frame}

\end{document}
